\section{数据来源与变量构造}

\subsection{数据来源与样本范围}

本文使用 CEPII BACI HS07 年度双边贸易数据 \citep{cepii_baci}，样本期为 2008--2024 年，进口国固定为中国。BACI 数据以 HS6 产品和双边贸易流为基本单位，适合从来源国、产品和年份三个维度刻画中国关键产品进口结构。本文保留年度口径，是为了保证 2008--2024 年样本期内不同产品和来源国之间具有可比性。

\subsection{产品选择与样本构造}

研究对象为四类半导体相关 HS6 产品，具体见表 \ref{tab:products}。这些产品覆盖半导体制造设备、处理器及控制器、存储器和其他电子集成电路，能够同时反映设备类与集成电路类产品的差异。

\begin{table}[H]
\centering
\caption{研究产品范围}
\label{tab:products}
\begin{tabular}{lll}
\toprule
HS6编码 & 产品类别 & 产品含义 \\
\midrule
848620 & 设备类 & 半导体器件或集成电路制造设备 \\
854231 & 集成电路类 & 处理器及控制器 \\
854232 & 集成电路类 & 存储器 \\
854239 & 集成电路类 & 其他电子集成电路 \\
\bottomrule
\end{tabular}
\end{table}


观察单元为出口国--产品--年份。对每个产品，本文保留样本期内曾向中国出口该产品的出口国，并补齐 2008--2024 年的年份维度；未出现的出口国--产品--年份贸易流按 0 处理。处理后的多产品平衡面板共 10608 行，其中正向贸易记录 5619 行，涵盖四个产品的所有历史出口国。

\subsection{变量定义与模型设定}

核心变量包括两类被解释变量：进口额对数 $\ln(Import_{i,p,t}+1)$ 用于衡量贸易规模变化，进口份额 $Share_{i,p,t}$ 用于反映来源国相对地位变化。核心解释变量为美国来源虚拟变量 $US_i$ 与三个政策节点虚拟变量（$Post2018_t$、$Post2022_t$、$Post2023_t$）的交互项，分别捕捉 2018 年实体清单扩容、2022 年出口管制新规和 2023 年持续收紧三个政策阶段对美国来源进口的相对影响。

控制变量包括出口国固定效应（吸收不随时间变化的国别特征）、产品固定效应（吸收产品层面的结构性差异）和年份固定效应（吸收宏观经济周期和全球半导体需求波动）。所有回归使用 HC1 异方差稳健标准误。
